% Source-aligned English for printed p.176.
\SourcePage{176}

\LayerTarget{tohoku:corollary:3.8.1:1:p2}
\noindent\textit{all the $U_{i_0\ldots i_p}$ are $F$-acyclic (that is, they
satisfy $H^p(U_{i_0\ldots i_p},F)=0$ for $p>0$).}

\LayerTarget{tohoku:def:point-indexed-cover-system}
We now restrict attention to coverings
$\mathcal U=(U_x)_{x\in X}$ indexed by the points $x$ of $X$, with
$x\in U_x$ for every $x$, and order them by writing
$\mathcal U\preceq\mathcal U'$ if $U'_x\subset U_x$ for every $x$.

\LayerTarget{tohoku:map:point-cover-refinement-complexes}
If $C_{\mathcal U}$ and $C_{\mathcal U'}$ are the corresponding
functor-complexes on $C^X$, there is then a functorial homomorphism
$C_{\mathcal U}\longrightarrow C_{\mathcal U'}$, hence a homomorphism of the
corresponding spectral functors. Immediate passage to the inductive limit gives

\smallskip
\noindent\textbf{Editorial correction.} The printed text gives
$U_x\subset U'_x$ while also writing
$C_{\mathcal U}\longrightarrow C_{\mathcal U'}$. Čech restriction maps have
this direction when $\mathcal U'$ refines $\mathcal U$, hence when
$U'_x\subset U_x$ for every $x$. That corrected order is used here and in the
following inductive limit.

\LayerTarget{tohoku:corollary:3.8.1:2}
\noindent\textsc{Corollary 2.} \textit{Let $X$ be any topological space.
There exists a spectral functor on the category $C^X$ of sheaves of abelian
groups on $X$, abutting to the graded functor $(H^n(X,F))$, whose initial term
is given by}

\LayerTarget{tohoku:eq:global-cech-spectral-E2}
\[
E_2^{p,q}(F)=\check H^p(X,H^q(F))
\]
\textit{(where $H^q(F)$ is the presheaf defined in Theorem 3.5.1). We have}

\LayerTarget{tohoku:eq:global-cech-spectral-E20q}
\[
E_2^{0,q}(F)=0\quad\textit{for }q>0.
\]

This last formula follows from the definition of $\check H^0(X,H^q(F))$ and
the following fact:

\LayerTarget{tohoku:lemma:3.8.2}
\noindent\textsc{Lemma 3.8.2.} --- \textit{Let $U$ be a neighborhood of $x$
and let $c^q\in H^q(U,F)$. Then there exists a neighborhood $V\subset U$ of
$x$ such that the image of $c^q$ in $H^q(V,F)$ is zero.}

\LayerTarget{tohoku:proof:3.8.2}
Indeed, take an injective resolution of $F\mid U$, let $C$ be the corresponding
complex, and represent $c^q$ by an element of $H^q(\Gamma(U,C))$ represented
by a cocycle $z\in\Gamma(U,C^q)$. Since $C$ is acyclic in degree $q$, the
restriction of $z$ to a suitable neighborhood $V$ of $x$ is a coboundary.
This proves the assertion, since $C\mid V$ is an injective resolution of
$F\mid V$ by Proposition 3.1.3, and hence
$H^q(V,F)=H^q(\Gamma(C\mid V))$.

The spectral sequence of Corollary 2 gives functorial homomorphisms

\LayerTarget{tohoku:map:global-cech-to-derived-cohomology}
\[
\check H^p(X,F)\longrightarrow H^p(X,F),
\]
and the formula

\LayerTarget{tohoku:eq:global-cech-E201}
\[
E_2^{0,1}=0
\]
shows that:

\LayerTarget{tohoku:corollary:3.8.1:3}
\noindent\textsc{Corollary 3.} \textit{The preceding homomorphisms are
bijective for $p=0$ or $1$ (as we already knew), and are monomorphisms for
$p=2$.}

\LayerTarget{tohoku:remark:paracompact-cech-comparison}
We also recover the fact that if $X$ is \textit{paracompact}, then
$\check H^p=H^p$; more precisely, \textit{the canonical homomorphisms above
are then isomorphisms.} Indeed, in this case one checks (using the fact that
the sheaf associated to the presheaf $H^q(F)$ is zero for $q>0$) that

\LayerTarget{tohoku:eq:paracompact-global-cech-E2-vanishing}
\[
E_2^{p,q}=0\quad\text{for }q>0;
\]
see [9].

\LayerTarget{tohoku:criterion:global-cech-edge-range}
Corollary 3 generalizes as follows: \textit{If
$\check H^p(X,H^q(F))=0$ for $0<q<n$, then the homomorphism
$\check H^i(X,F)\longrightarrow H^i(X,F)$ is an isomorphism for $i\leq n$ and
a monomorphism for $i=n+1$.} We conclude (with H. Cartan):

\LayerTarget{tohoku:corollary:3.8.1:4:p1}
\noindent\textsc{Corollary 4.} \textit{Let $\mathcal U$ be a collection of
open sets forming a basis for the topology of $X$, and let $F$ be a sheaf of
abelian groups on $X$ such that for every nonempty sequence}

\TohokuContinuation{177}
